% =========================================================
\chapter{C07 --- Introduction aux diodes}
% =========================================================

\begin{rembox}[title=En bref]
\textbf{De quoi parle ce chapitre ?} Les diodes ouvrent les \textbf{dispositifs à
semi-conducteur} : composants passifs (qui ne font que dissiper de l'énergie) et actifs (les
transistors, qui amplifient). La méthode vue ici (signaux continus + alternatifs) resservira
pour les transistors. \\
\textbf{Les trois objectifs} :
\begin{enumerate}
  \item \textbf{un composant non linéaire} : \og bienvenue dans un monde sans loi d'Ohm
        linéaire\fg{} ; la diode suit une loi \textbf{exponentielle} $I_D = I_S e^{U_D/nU_T}$ ;
  \item \textbf{la modélisation analytique} : avec une exponentielle, la résolution directe
        devient impossible, même pour une diode et une résistance ; la \textbf{droite de
        charge} donne l'intuition, puis des \textbf{modèles} simples ;
  \item \textbf{la dualité standard vs Zener} : la Zener limite le passage du courant dans
        les \textbf{deux} sens.
\end{enumerate}
\end{rembox}

\section{La diode normale}

\slide{c07-03}{Symbole et caractéristique : sens direct, sens inverse, zone Zener.}

\begin{defnbox}[title=Symbole et caractéristique]
Le symbole est une \textbf{flèche} qui indique le sens naturel du courant : l'entrée est
l'\textbf{anode}, la sortie la \textbf{cathode}. Le constructeur donne aussi la
\textbf{puissance maximale} $P_{MAX}$ ($P = UI$) : quel que soit le type de diode, un courant
trop élevé la détruit. Trois zones de comportement :
\begin{enumerate}
  \item vers \textbf{0,7\,V} (diode au silicium) : le courant grimpe très vite pour le
        moindre accroissement de tension (\textbf{sens direct}) ;
  \item entre une tension négative limite et 0,7\,V : quasiment pas de courant (seulement
        $I_S$, insignifiant) : la diode est \textbf{bloquée} ;
  \item au seuil négatif $-V_Z$ (\textbf{tension Zener}) : le courant redevient très
        important, en sens inverse (\textbf{effet Zener ou avalanche}, effet tunnel\ldots).
        Destructeur ou non selon la construction ; s'il est prévisible et non destructif
        (réversible), il est exploitable : ce sont les \textbf{diodes Zener}.
\end{enumerate}
\end{defnbox}

\slide[0.7]{c07-04}{Parenthèse semi-conducteurs (pour en savoir plus).}

